\section{Methodology}

\subsection{Simulation Environment and Network}

All experiments run in SUMO~\cite{lopez2018sumo}; \CO{} comes from its HBEFA3
implementation~\cite{krajzewicz2015emissions}. Policies act through the
simulator's online control interface, so every routing decision is taken inside
the running simulation from information available at that instant.

Figure~\ref{fig:network} shows the environment: one origin--destination pair
served by three structurally different corridors between a common diverge and
merge, crossed by four minor streets carrying independent demand through the
same signals, so a policy that loads a corridor also delays traffic it does not
route. Because every junction permits straight-through movements only, the
loopless path set is exactly these three --- the candidate catalogue is complete,
not a $k$-shortest approximation.

Two structural choices were fixed before the experiment, both identified by
capacity validation. The destination portal gives each approach a dedicated exit
lane, so an alternative corridor is not penalised by junction priority rather
than by its own capacity; and the diverge carries no turn-radius speed ceiling,
which in schematic node geometry would cap the bypass at a value belonging to
the drawing rather than to the road.

\paragraph{Measured capacity.} Each corridor was driven to saturation alone and
its discharge measured: \NumSatArterial~veh/h/lane on the arterial and
\NumSatNorth~veh/h/lane on the slower street, consistent to within
\NumSatSpread{} across green ratios, and \NumCapBypass~veh/h on the
uninterrupted bypass --- below queue-discharge saturation flow, as free-flow
headways require. The demand levels of Table~\ref{tab:factors} follow from these
measurements rather than from a textbook constant.

\begin{figure}[tbp]
\centering\includegraphics[width=\textwidth]{fig1_network.pdf}
\caption{The scenario network. One origin--destination pair is served by three
structurally different corridors between a common diverge and merge, crossed by
four minor streets carrying independent demand through the same signals.
Geometry, lane counts and speed limits are read from the built network file;
corridor capacities are measured rather than assumed.}
\label{fig:network}
\end{figure}

\subsection{Routing Policy Portfolio}

Table~\ref{tab:portfolio} gives each policy's objective, information
requirement, intended strength and expected failure mode. \emph{The failure
modes were specified before the experiment}, and each factor exists to make one
of them reachable: P1 $\rightarrow$ capacity shortage on the shortest corridor;
P2 $\rightarrow$ information lag and herding; P3 $\rightarrow$ unevenly
distributed alternative capacity; P4 $\rightarrow$ travel-time variability under
disruption.

All four read traffic state from one shared buffer: a policy acting at time $t$
may read link travel times sampled no later than $t-\Delta$ over a 120\,s
window, falling back to free-flow values before any admissible sample exists. No
policy sees the future or any realised outcome of the run it is part of. Policy
parameters were frozen before the experiment ($\varepsilon=\NumEpsFrozen$ and a
120\,s assignment window for P3, $\lambda=1.0$ for P4).

Each policy was validated before use by twelve known-answer tests covering the
observation lag, the free-flow fallback, and each policy's defining behaviour.
A two-path integration test inside the simulator, with equal path lengths and a
2:1 capacity ratio, confirms that P3 splits demand \NumTwoPathRatio{} against
that ratio while P1 commits the entire cohort to one path.

\begin{table}[t]
\caption{The routing-policy portfolio. Every policy is an established family;
none is introduced here. The last column states a hypothesis fixed before the
experiment.}
\label{tab:portfolio}
\centering\small
\begin{tabularx}{\textwidth}{@{}l l l X X@{}}
\toprule
& Objective & Information used & Intended strength & Hypothesised failure mode\\
\midrule
P1 & $\min_p L(p)$ & none &
  needs no infrastructure and cannot oscillate &
  loads the shortest corridor past its capacity\\[2pt]
P2 & $\min_p \hat\mu_p(t-\Delta)$ & lagged link travel times &
  follows congestion as it forms &
  moves the whole guided cohort together on stale information\\[2pt]
P3 & $\min_{p\in\mathcal{A}} \max_{e\in p}\rho_e$ &
  lagged travel times and flows, plus its own recent assignments &
  spreads load when capacity is unevenly distributed &
  diverts to longer corridors when no capacity shortage exists\\[2pt]
P4 & $\min_p [\hat\mu_p + \lambda\hat\sigma_p]$ & lagged link travel times &
  avoids corridors whose travel time is volatile &
  pays a detour for variance that carries no risk\\
\bottomrule
\end{tabularx}
\end{table}

\subsection{Experimental Factors and Common Random Numbers}
\label{sec:factors}

Six factors (Table~\ref{tab:factors}) are crossed fully: \NumContexts{}
contexts, \NumSeeds{} seeds, four policies, \NumRunsMain{} runs. The demand
levels span uncongested to oversaturated conditions relative to the measured
capacity (Figure~\ref{fig:design}).

The vehicle set, departure times, guided/unguided labelling, habitual corridor
choice, background traffic and the incident realisation depend on the seed and
context alone, never on the policy, so two runs of a context under different
policies face identical vehicles and every comparison is paired. Unguided
drivers follow a habitual corridor drawn once from a frozen logit over free-flow
path times; at penetration $p$, a fraction $1-p$ of corridor demand is therefore
inert with respect to the policy under test.

\paragraph{Disruption as risk, not as a known event.} When the incident regime
is active, one lane closure is drawn --- location, onset and duration --- from a
declared distribution seeded by the run seed. Within a context and seed the
realisation is identical across policies, preserving the pairing; across seeds
it differs, so the context carries travel-time risk no policy can know in
advance. An incident occurring identically in every replication would be a known
event rather than a risk, and could not test a reliability-aware policy.

\paragraph{Statistical treatment.} The experimental unit is the context. An
advantage of $a$ over $b$ is \emph{resolved} only when the paired seed
differences satisfy $|\overline{d}| > 2\,\mathrm{SE}(\overline{d})$; unresolved
differences are treated as ties, so no claim rests on a difference the seeds
cannot separate.

\begin{table}[t]
\caption{Scenario factors and levels. Demand levels are expressed as a fraction
of the measured network capacity, which ranges from 0.31 to 1.10 at the
median capacity configuration. Full factorial: 648 contexts.}
\label{tab:factors}
\centering\footnotesize
\begin{tabular}{@{}l l l@{}}
\toprule
Factor & Levels & Mechanism activated\\
\midrule
Demand (veh/h) & 1200, 1800, 2400, 3000, 3600, 4200 & capacity shortage\\
Green ratio $g/C$ & 0.35, 0.50, 0.65 & signalised capacity\\
Penetration & 0.2, 0.5, 0.8 & cohort size; herding\\
Information lag $\Delta$ (s) & 30, 120, 300 & staleness; oscillation\\
Incident regime & none, arterial & travel-time risk\\
Alt.\ capacity & 1, 2 lanes & load to balance\\
\bottomrule
\end{tabular}
\end{table}

\begin{figure}[tbp]
\centering\includegraphics[width=\textwidth]{fig3_design.pdf}
\caption{Demand levels against measured network capacity. The shaded band spans
the six capacity configurations produced by the green-ratio and
alternative-capacity factors; the line marks the median configuration. The grid
spans uncongested to oversaturated conditions, and the levels follow from
measurement rather than from a chosen range.}
\label{fig:design}
\end{figure}

\subsection{Outcome Criteria and Pareto Evaluation}
\label{sec:criteria}

Three criteria (Table~\ref{tab:criteria}) capture different burdens.
\textbf{C1 mean journey time} is the travel burden per vehicle, including
insertion delay --- a policy that oversaturates the network produces vehicles
that cannot enter it, and excluding their wait would credit the policy for the
queue it caused. \textbf{C2 \CO{}} is the environmental outcome. \textbf{C3
total stopped delay} is the stationary burden, time below 0.1\,m/s in queues,
distinct from time spent travelling slowly.

All three are measured over every vehicle in the network, guided or not, with
intended departure in the measurement window, so a policy cannot be rewarded for
helping guided vehicles at the expense of background traffic. Policies are
compared by \textbf{Pareto dominance first}; the scalar cost used for regret is
C1, which requires no weights. Scalarisation appears only as sensitivity,
through three declared normalised preference profiles. No monetary value of
time, value of reliability or carbon price is used: none was available from a
verified source, and inventing one would make the conclusion an artefact of the
invented number.

\begin{table}[t]
\caption{The three decision criteria. All are measured over the same cohort:
every vehicle in the network, guided or not, whose intended departure falls in
the measurement window.}
\label{tab:criteria}
\centering\footnotesize
\begin{tabularx}{\textwidth}{@{}l l X@{}}
\toprule
& Criterion & Measurement\\
\midrule
C1 & Mean journey time (s/veh) &
  mean over the cohort of insertion delay plus in-network time. Insertion delay
  is included: a policy that oversaturates the network produces vehicles that
  cannot enter it, and excluding their wait would credit the policy for the
  queue it caused.\\[2pt]
C2 & \CO{} emissions (kg) &
  total \CO{} over the cohort, HBEFA3 passenger-car petrol Euro-4 class.\\[2pt]
C3 & Total stopped delay (veh\,s) &
  total time spent below 0.1\,m/s --- the time vehicles spend \emph{stationary}
  in queues, as distinct from the time they spend travelling slowly.\\
\bottomrule
\end{tabularx}
\end{table}

\subsection{Mechanistic Rule}

Each context is expressed through eight dimensionless descriptors computable
before activation: network saturation $x_1$; shortest-corridor saturation
$x_2 = D/\mathrm{cap}_C$; alternative capacity share $x_3$; effective green
ratio $x_4$; penetration $x_5$; information staleness
$x_6 = \Delta/T^{\mathrm{ff}}_C$; incident regime $x_7$; and controlled-flow
saturation $x_8 = p\,D/\mathrm{cap}_C$, the share of the shortest corridor's
capacity the guided cohort alone would consume. Capacities are the measured
values, so $x_1$, $x_2$ and $x_8$ are genuine degrees of saturation. A
route-overlap index is omitted: the corridors share only access and egress
edges, so overlap is structurally constant here.

The mechanistic rule B1 is an axis-aligned region rule on $x_1$--$x_8$ of depth
at most two, fitted on training contexts only. Thresholds are reported as
\emph{brackets between adjacent sampled levels}, since a grid cannot locate a
boundary more finely than its own spacing.

\subsection{Context-Aware Selector}

The learned selector predicts \emph{advantage}, not identity:
\begin{equation}
\Delta_a(x) \;=\; C(x, a_{\mathrm{SBS}}) - C(x,a),
\qquad \hat{\pi}(x) = \arg\max_a \hat{\Delta}_a(x),
\label{eq:advantage}
\end{equation}
with $a_{\mathrm{SBS}}$ determined on training contexts only. A positive
$\Delta_a$ means $a$ beats the fixed policy a deployed system would otherwise
run, so the learned quantity is the one a decision-maker needs. The model is
gradient-boosted decision trees~\cite{friedman2001gbm,ke2017lightgbm}, one
regressor per policy, with hyperparameters fixed before evaluation and never
tuned on a test fold. With a few hundred context-level instances and full
feedback from every policy, higher capacity is neither needed nor defensible; we
use no graph neural network, deep network, reinforcement learning or contextual
bandit.

The ladder is B0 fixed policy, B1 mechanistic rule, B2 depth-3 tree, B3
multinomial logistic regression, B4 the advantage selector and B5 its selective
variant, with VBS as reference. \textbf{The primary comparison is B4 against
B1}: B0 is a low bar, and beating it shows only that the decision varies with
context at all.

\subsection{Validation, Uncertainty, and Distribution Shift}

Evaluation is leave-one-context-out. Inside each fold, standardisation, the
identity of the SBS, tree thresholds, conformal quantiles and the support
envelope are recomputed from training data alone. No seed of a held-out context
appears in training, and no realised outcome is used as a feature.

\paragraph{Two gates.} B5 acts only if both pass. The \emph{support gate}
measures mean distance to the five nearest training contexts in standardised
descriptor space and abstains above the 95th percentile of the training fold's
own distances; it is a feature-space extrapolation detector. The
\emph{confidence gate} is a split-conformal interval on the predicted advantage
($\alpha=0.10$, calibrated on 30\% of the training fold); the selector acts only
when the lower bound on the advantage over the fixed policy is positive, and it
is calibrated under exchangeability, hence in-distribution. If either gate
fails, the fixed policy is used. Neither gate detects concept shift, and the
conformal guarantee does not survive covariate
shift~\cite{shimodaira2000,tibshirani2019covshift}, so we report gate behaviour
rather than coverage guarantees outside the training distribution.

\paragraph{Shift splits.} Eight held-out splits are evaluated separately and
never pooled: demand above and below the training range, an interior demand
level as an interpolation control, unseen information lag, unseen penetration,
unseen green ratio, training without disruption and testing under it, and a
disruption on a corridor that carries none in training. Pooling them would
average an interpolation control with a domain shift and hide the only thing the
experiment can say, which is that they differ.
